% =====================================================================
% TmFeO3 2DCS — main text (letter format)
% Long-form derivations: ../tmfeo3_foundation.tex; supplement: supp.tex
% =====================================================================
\documentclass[aps,prl,twocolumn,superscriptaddress,longbibliography]{revtex4-2}

\usepackage{amsmath,amssymb,bm}
\usepackage{graphicx}
\usepackage{xcolor}
\usepackage{hyperref}
\hypersetup{colorlinks=true,linkcolor=blue,citecolor=blue,urlcolor=blue}

\newcommand{\Fe}{\mathrm{Fe}}
\newcommand{\Tm}{\mathrm{Tm}}
\newcommand{\ket}[1]{|#1\rangle}
\newcommand{\ve}[1]{\boldsymbol{#1}}
\newcommand{\qafm}{q_{\rm AFM}}
\newcommand{\qfm}{q_{\rm FM}}
\newcommand{\Mnl}{M_{\rm NL}}

\begin{document}

\title{Reading a magnet's Hamiltonian term by term:\\
two-dimensional terahertz spectroscopy as vertex spectroscopy of
TmFeO$_3$}

\author{Working draft}
\affiliation{---}

\date{\today}

\begin{abstract}
One-dimensional spectroscopy measures the \emph{spectrum} of a
Hamiltonian; two-dimensional coherent spectroscopy (2DCS) measures its
\emph{vertices}---the individual nonlinear terms that convert one
coherence into another.  We show that in the rare-earth orthoferrite
TmFeO$_3$ this correspondence can be made one-to-one.  Time-reversal
parity and space-group irrep bookkeeping reduce every symmetry-allowed
Fe--Tm coupling to exactly two conversion vertices, a static two-magnon
exchange $WS^2\lambda$ and a field-gated $\kappa^B BS\lambda$; the site
mirror of the Tm $4c$ position then assigns each light polarisation to
one block of the crystal-field qutrit, fixing what each channel can
excite and what it can radiate.  Driven by the measured single-cycle
THz waveform, with every damping rate set by a measured linewidth and
every emission weight by a measured oscillator strength, one
Hamiltonian reproduces the cross-polarised inventories of both field
geometries without retuning, and---with a detection operator whose
three terms are each required by symmetry---four of the six
same-polarised amplitudes to within $15\%$.  Each peak identifies a
vertex and each structural absence a selection rule, while the residual
discrepancy reduces to a single dynamical ratio.  We give the
falsifiable signature of every assignment.
\end{abstract}

\maketitle

\emph{Introduction.}---Linear response cannot distinguish two
Hamiltonians with the same normal modes: absorption measures pole
positions and dipole weights, and any coupling that merely hybridises
modes is absorbed into the mode basis.  What it can never see is a
\emph{conversion}: a coherence prepared on one transition and
re-emitted on another.  That is precisely what a cross peak in
two-dimensional coherent spectroscopy certifies~\cite{Lu2017,Mukamel},
and a conversion requires a specific anharmonic vertex.  A cross peak
at $(\omega_\tau,\omega_t)$ is therefore not a feature to be fitted but
a \emph{measurement of one coupling term}, with the delay axis naming
the stored coherence, the detection axis the emitter, and the signed
side constraining the pathway order.  If symmetry makes the space of
candidate vertices finite---and in a crystal it does---a
polarisation-resolved 2D atlas can be read as the Hamiltonian itself.

TmFeO$_3$ is an unusually clean testbed because two qualitatively
different subsystems share one THz octave.  Below the spin reorientation
the Fe sublattice orders in $\Gamma_2$
$(F_x,C_y,G_z)$~\cite{Yamaguchi,Kimel2004,Zhang2016} with
quasi-ferromagnetic and quasi-antiferromagnetic magnons at
$\qfm=0.38$ and $\qafm=0.90$\,THz, while the three lowest Tm$^{3+}$
crystal-field (CEF) singlets form a qutrit with transitions
$E_{12}=0.50$, $E_{23}=0.70$, $E_{13}=1.20$\,THz.  A single-cycle THz
pulse covers all five modes, and measurements in the cross- and
same-polarised channels of two field geometries return distinct,
reproducible inventories.

\emph{Model.}---The Fe$^{3+}$ ($S=5/2$) sublattice is classical, with
the calibrated orthoferrite Hamiltonian (exchange, single-ion
anisotropy, Dzyaloshinskii--Moriya) reproducing $\qfm$, $\qafm$ and
$\Gamma_2$.  Each Tm$^{3+}$ is truncated to its three lowest singlets
and carried as the Gell-Mann vector
$\lambda^a=\langle\Lambda^a\rangle$.  Because every single-ion term is
linear in the $\Lambda^a$, the classical $\vec\lambda$ equation of
motion \emph{is} the exact von Neumann equation of the driven, damped
qutrit, and a Boltzmann ensemble over level seeds is the exact thermal
density matrix [Supplemental Material (SM), Sec.~S1].  The Fe--Tm
sector is obtained by projecting the physical exchange
$\ve S^{T}\!A\,\ve J$ into the qutrit in a transported gauge over the
four Tm sublattices [SM, Secs.~S2--S3].

\emph{Two vertices survive.}---A conversion vertex must reach the
bright coordinate, be sourced by the pumped mode, be time-reversal
even, be $\Gamma_1$-allowed in the ordered state, carry the frequency
mismatch, and engage both pulses to survive the subtraction
$\Mnl=M_{AB}-M_A-M_B$.  These six requirements exhaust the coupling
space [SM, Sec.~S4].  The static bilinear $K^-S\lambda$ fails
threefold; the quadratic families $W_4$, $W_1^{xy}$, $W_6$,
$\kappa^E$ are exactly inert.  What remains is
\begin{align}
  H_W&=W^{yy}_1\,S_y^2\,\lambda^1, &
  H_{\kappa}&=\kappa^B_{1y}\,B_y(t)\,S_y\,\lambda^1 ,
  \label{eq:vertices}
\end{align}
and the data require both.

\emph{Symmetry assigns each polarisation a block.}---The Tm $4c$ site
has point symmetry $m$.  The projected moment matrix---an output of the
construction, not an input---has $\lambda^2$ carrying only $\mu_z$ and
$\lambda^{5,7}$ only $\mu_{x,y}$, which forces $\ket1,\ket2$ to share a
mirror parity and $\ket3$ to be opposite.  Every coupling follows:
$1\!\leftrightarrow\!2$ is reached by $m_z$ and $d_{x}$, hence by the
$(E\!\parallel\!a,H\!\parallel\!c)$ polarisation;
$1\!\leftrightarrow\!3$ and $2\!\leftrightarrow\!3$ by $m_{x}$ and
$d_z$, hence by $(E\!\parallel\!c,H\!\parallel\!a)$.  Each polarisation
therefore owns one block and---by reciprocity---drives and reads the
same one, so cross-polarised detection reads exactly the block the
pulse did not drive.  This is why cross-polarisation isolates
conversion, and it independently predicts the subalgebra closure we
observe in geometry B, where $\lambda^{4\text{--}7}$ vanish to machine
precision.

All spectra use the \emph{measured} pump (spectrum peaked at
$0.53$\,THz, no DC), each coherence damped at its measured linewidth,
and each emitter weighted by its measured oscillator strength
[SM, Sec.~S5].  Peaks are identified blindly as 2D local maxima with
prominence $\geq1.5$; the $\tau$-averaged signal is subtracted at each
detection time, removing the $\omega_\tau=0$ row.  One fluence serves
every geometry, fixed by requiring the cross channel to be perturbative
[SM, Sec.~S10].

\begin{figure*}[t]
\includegraphics[width=\textwidth]{figs/atlas_final.png}
\caption{The atlas: four polarisation channels from one Hamiltonian at
one operating point.  Detection is resolved by the emitted magnetic
polarisation, so each panel reads one moment component [SM, Tab.~S3].
Top left: geometry-A cross ($m_z$)---the magnetoelectric
$(\qafm,E_{12})$ peak of $W^{yy}_1$ and, at parity, the
$(E_{12},2E_{12})$ second harmonic of the non-centrosymmetric Tm site;
Fe is radiatively silent here ($F_z\equiv0$).  Top right: geometry-B
cross ($m_x$)---the magnon--magnon $(\qfm,\qafm)$ peak and the
$\kappa^B$-gated $(E_{12},\qafm)$ transfer.  Bottom left: geometry-A
same-polarised.  Bottom right: the geometry-B same-polarised
prediction, a single harmonic peak whose strength measures the Tm
hyperpolarisability.  Dashed guides mark the five mode energies;
crosses are blind-census maxima.}
\label{fig:atlas}
\end{figure*}

\emph{The cross channels are the anchors.}---In geometry A the
Zeeman-launched $\qafm$ magnon precesses through the delay and is
transcribed onto the bright Tm coordinate by $W^{yy}_1$, whose
two-pulse cross term rectifies to
$\Mnl\propto\cos(\qafm\tau)[1-\cos E_{12}t]$: a factorised, rank-one
peak at $(\qafm,E_{12})$.  The magnon never radiates here---it has no
$c$-axis moment---so this is a magnetoelectric object, an E1-dark
magnon made visible by borrowing the Tm dipole, and its existence
measures $W^{yy}_1$.  The $E_{12}$ quantum comes from the pulse edge:
stretching the pulse fourfold at fixed amplitude collapses the peak by
$4.5$ decades, while switching off the direct Tm drive changes it by
$<10^{-4}$ [SM, Sec.~S6].  The second reported feature sits at
$2E_{12}=1.0000$\,THz and is second-harmonic emission: a linear moment
cannot radiate a harmonic, but the Tm $4c$ site has no inversion (Fe at
$4b$ does), so $m_z=\mu\lambda^2+\beta\lambda^1\lambda^2$ is allowed,
is $\mathcal T$-odd as a magnetisation must be, and radiates at
$2\omega_{12}$.  It contributes essentially one peak, and
$\beta\simeq7\mu$ from the observed parity gives
$(\qafm,E_{12})=1.00$ and $(E_{12},2E_{12})=0.97$.

In geometry B both peaks emit at $\qafm$ through the magnon dipole.
$(\qfm,\qafm)=1.00$ is a pure two-magnon conversion and measures the Fe
sector alone.  $(E_{12},\qafm)=0.45$ isolates the gated vertex: with
$\kappa^B_{1y}$ alone it is present, with the static hybridisation
alone it is absent, and \emph{no} strength of the latter recovers
it---a sweep shows a static vertex only merges and shifts the delay
rows, because a mixer moves peaks while a gated converter adds them at
unshifted lines [SM, Sec.~S7].  The same $B_z$ gate makes $\kappa^B$
identically silent in geometry A.  Both cross inventories are
reproduced without retuning and, importantly for what follows, are
unchanged under every drive variant we have tested, because they are
fed by conversion rather than by the direct Tm drive.

\emph{The same-polarised channel.}---Here the analyser reads $m_x$,
which receives the Tm $a$-axis moments and the Fe $F_x$.  Three
contributions are each required by symmetry, and each carries one
measured feature: the CEF moment ($E_{23}$ and $E_{13}$ emission), the
Fe $F_x$ (magnon emission), and---because $E\!\parallel\!c$ drives
$\lambda^2$ through the exchange-induced $d_z^{\rm eff}$ that the
$\Gamma_2$ order permits below $T_N$---an $E_{12}$ emission term
demanded by reciprocity.  That last term is not a device: the measured
map contains a peak at $(\qafm,E_{12})$, i.e.\ $E_{12}$ emission in
this channel, at $0.30$ of maximum.  Fitting the three weights to the
digitised measurement with the cross channels held fixed gives
\begin{center}\footnotesize
\renewcommand{\arraystretch}{1.05}
\begin{tabular}{lcc}
\hline
peak & experiment & model\\
\hline
$(E_{12},E_{23})$ & $1.00$ & $1.00$\\
$(E_{12},\qfm)$ & $0.45$ & $0.38$\\
$(E_{12},\qafm)$ & $0.40$ & $0.46$\\
$(\qafm,E_{12})$ & $0.30$ & $0.32$\\
$(\qafm,E_{13})$ & $0.30$ & $0.09$\\
$(E_{12},E_{13})$ & $0.15$ & $0.44$\\
\hline
\end{tabular}
\end{center}
Four of six lie within $15\%$, and the fit drives the residual
$\mu^z_{13}$ admixture to zero---the $1\!\to\!3$ transition is
$c$-axis dark, consistent with the complete absence of
$\omega_\tau=E_{13}$ rows from the measured map, which a bright
$\mu^z_{13}$ would have produced at first order.

\emph{What remains.}---The two failures are not independent: both
peaks are detected at $E_{13}$ and therefore share an emission weight,
so their ratio is set by the dynamics alone,
\begin{equation}
  \frac{(E_{12},E_{13})}{(\qafm,E_{13})}:\qquad
  \text{model }4.9,\qquad \text{experiment }0.5 .
  \label{eq:gap}
\end{equation}
No detection weight can move it.  The model routes too much $E_{13}$
emission through the $E_{12}$-excitation row relative to the magnon
row, and closing that factor of ten is the one open item of the atlas.

\emph{The atlas as an overdetermined measurement.}---Fifteen
independent observations constrain five couplings
($W^{yy}_1$, $W^{xz}_1$, $\kappa^B_{1y}$, $W_3$, $\beta$), one
selection rule and two drive parameters; SM Sec.~S11 gives a
peak-by-peak account.  The falsifiable content: any peak emitting at
$\qafm$ inherits the $15\times$ sharper magnon linewidth;
$(\qfm,\qafm)$ is temperature-flat while $(E_{12},\qafm)$ falls as
$n_1-n_2$; rotating either polarisation away from $c$ must revive the
$\omega_\tau=E_{13}$ rows; the $2E_{12}$ line tracks $E_{12}$ and so
does \emph{not} soften at the spin reorientation, whereas every feature
relying on $d_z^{\rm eff}\propto\langle F_x\rangle$ must vanish there;
the geometry-B same-polarised channel should be a single harmonic peak
whose strength measures $\beta$; and adiabatically long pulses collapse
the geometry-A cross peak.

\emph{Conclusion.}---Two symmetry-selected Fe--Tm vertices, one
Fe-sector magnon--magnon coupling, a mirror-parity rule assigning each
polarisation to one qutrit block, and a hyperpolarisability allowed by
the absence of inversion at the rare-earth site account for the
measured 2DCS atlas of TmFeO$_3$---computed from measured inputs, with
no free lineshape parameters and one residual dynamical ratio.  The
program---bound the vertex space by symmetry, assign every peak and
every absence, and let the atlas overdetermine the model---applies to
any magnet whose 2D map spans distinct subsystems, and turns 2DCS from
a spectroscopy of modes into a spectroscopy of Hamiltonian terms.

\begin{acknowledgments}
Derivations, numerical protocol, the peak-by-peak account and the
complete falsification ledger are in the Supplemental Material; the
long-form note and all run recipes accompany the data.
\end{acknowledgments}

\begin{thebibliography}{9}
\bibitem{Lu2017} J. Lu, X. Li, H. Y. Hwang, B. K. Ofori-Okai,
  T. Kurihara, T. Suemoto, and K. A. Nelson,
  Phys.\ Rev.\ Lett.\ \textbf{118}, 207204 (2017).
\bibitem{Mukamel} S. Mukamel, \emph{Principles of Nonlinear Optical
  Spectroscopy} (Oxford University Press, New York, 1995).
\bibitem{Yamaguchi} T. Yamaguchi,
  J.\ Phys.\ Chem.\ Solids \textbf{35}, 479 (1974).
\bibitem{Kimel2004} A. V. Kimel, A. Kirilyuk, A. Tsvetkov, R. V.
  Pisarev, and Th. Rasing, Nature \textbf{429}, 850 (2004).
\bibitem{Zhang2016} \emph{Resolving the spin reorientation and
  crystal-field transitions in TmFeO$_3$ with terahertz transient},
  Sci.\ Rep.\ \textbf{6}, 23648 (2016).
\bibitem{MP2018} \emph{Terahertz magnon-polaritons in TmFeO$_3$},
  ACS Photonics \textbf{5}, 1375 (2018).
\end{thebibliography}

\end{document}
